\documentclass[a4paper,12pt]{article}

% --- Codificação e Idioma ---
\usepackage[utf8]{inputenc}
\usepackage[T1]{fontenc}
\usepackage[brazil]{babel} % Essencial para hifenização correta em PT-BR
\usepackage{helvet}
\renewcommand{\familydefault}{\sfdefault} % Fonte Arial-like
\usepackage{indentfirst}

% --- Matemática e Unidades ---
\usepackage{amsmath,amsfonts}
\usepackage{siunitx}
\sisetup{
    locale = FR,                 % Usa vírgula decimal
    per-mode = symbol,           % Barra para divisão de unidades
    separate-uncertainty = true, % Incerteza +/-
    detect-all                   % Fonte igual ao texto
}

% --- Tabelas e Gráficos ---
\usepackage{geometry}
\geometry{margin=2.5cm}
\usepackage{graphicx}
\usepackage{booktabs} % Tabelas profissionais
\usepackage{enumitem}
\usepackage{tikz}
\usetikzlibrary{calc, positioning, arrows.meta}

% --- Links e Metadados do PDF ---
\usepackage{hyperref}
\hypersetup{
    % Configuração para aceitar acentos nos metadados
    pdfencoding=auto, 
    unicode=true,
    % Metadados do arquivo PDF
    pdftitle={Desenvolvimento e Metodologia do Reômetro Capilar v4.1.1},
    pdfauthor={Bruno Kenji Nishitani Egami},
    pdfsubject={Reologia Capilar e Processamento de Dados},
    pdfkeywords={Reologia, Cerâmica, Python, Instrumentação, DIY, LM4040},
    % Configuração visual dos links
    colorlinks=true,
    linkcolor=black,
    citecolor=blue,
    filecolor=magenta,
    urlcolor=blue
}

% --- Cabeçalho e Dados de Contato ---
\title{Desenvolvimento e Metodologia do Reômetro Capilar v4.1.1}
\author{
    \textbf{Bruno Kenji Nishitani Egami} \\
    \small \href{mailto:bruno.egami@farroupilha.ifrs.edu.br}{bruno.egami@farroupilha.ifrs.edu.br} \\
    \small \href{https://github.com/bruno-egami}{GitHub: @bruno-egami} $\cdot$ 
    \href{https://www.instagram.com/bruno.egami/}{Instagram: @bruno.egami}
}
\date{Abril de 2026}

% ---------------------------------------------------------------------
\begin{document}

\maketitle

\noindent
\href{https://github.com/bruno-egami/Reometro_Capilar/tree/Reometro-capilar-Transdutor_Pasta}{\textbf{Acesse aqui o repositório específico do Projeto}}

\section*{1. Introdução}
A caracterização reológica de pastas cerâmicas é essencial para entender seu comportamento durante processos de conformação, como a extrusão. O presente relatório descreve as equações utilizadas em um conjunto de scripts Python desenvolvidos para processar os dados experimentais obtidos com um reômetro capilar (\textit{DIY}) de baixo custo.

O objetivo é converter medidas diretas (como massa extrudada e pressão de extrusão) em parâmetros reológicos interpretáveis (como viscosidade e taxa de cisalhamento).

\subsection*{1.1 Aplicações em Manufatura Aditiva}
O controle reológico preciso é fundamental para tecnologias emergentes como a impressão 3D cerâmica, especialmente na técnica \textit{Direct Ink Writing} (DIW) ou \textit{Robocasting}. Esta técnica requer pastas com comportamento pseudoplástico que fluam sob cisalhamento (extrusão pelo bico), mas recuperem rapidamente sua viscosidade após a deposição para suportar as camadas subsequentes (\textit{buildability}).

A caracterização sistemática do comportamento reológico da matéria-prima base é, portanto, um pré-requisito para o desenvolvimento de formulações otimizadas para manufatura aditiva.

\subsection*{1.2 Histórico de Desenvolvimento do Equipamento}
O reômetro capilar utilizado neste estudo é resultado de um processo evolutivo de três gerações:

\begin{itemize}
  \item \textbf{V1 (Protótipo Inicial):} A primeira versão utilizava medição indireta de pressão através de uma célula de carga (extensômetros e módulo HX711). Esta configuração permitia estimar a força aplicada durante a extrusão. O sistema foi validado experimentalmente demonstrando boa concordância com reômetros comerciais.
  
  \item \textbf{V2 (Transdutor Único):} Na segunda geração, a célula de carga foi substituída por um transdutor de pressão de linha (faixa de \SI{0}{psi} a \SI{150}{psi}) para leitura direta da pressão do ar comprimido. Esta modificação simplificou o sistema e melhorou a precisão das medições.
  
  \item \textbf{V3 (Atual - Duplo Transdutor e Referência AREF):} A versão atual incorpora um segundo transdutor posicionado diretamente na câmara da pasta (sensor "Pasta"), além do transdutor de linha (sensor "Linha"). Para mitigar a instabilidade da tensão USB (que flutua entre \SI{4,75}{V} e \SI{5,25}{V}), foi integrada uma referência de tensão externa de precisão \textbf{LM4040-4.096} conectada ao pino AREF do Arduino, garantindo que o ADC de 10 bits opere com uma base estável de \SI{4,096}{V} independentemente da fonte de alimentação. Esta configuração elimina erros sistemáticos de leitura e possibilita a aplicação rigorosa das correções de Bagley e Mooney.
\end{itemize}

\begin{figure}[htbp]
\centering
\includegraphics[width=0.85\textwidth]{esquema_reometro_v4.png}
\caption{Esquema de funcionamento do reômetro capilar, mostrando a linha pneumática, os sensores de pressão (Linha e Pasta) e o sistema de aquisição digital.}
\label{fig:esquema_reometro}
\end{figure}

O controle do ensaio e o processamento de dados são realizados por um conjunto de scripts em Python.

\subsection*{1.2.1 Processamento de Sinais e Filtragem Digital}
Os dados analógicos provenientes dos transdutores são digitalizados pelo ADC do Arduino a uma frequência de aproximadamente \SI{10}{Hz}. Para reduzir ruídos elétricos e oscilações pneumáticas de baixa magnitude, é aplicado um filtro de Média Móvel Exponencial (EMA) diretamente no firmware:

\begin{equation}
    S_t = \alpha \cdot Y_t + (1 - \alpha) \cdot S_{t-1}
\end{equation}

\textbf{Onde:}
\begin{itemize}[leftmargin=1.5em]
  \item $S_t$ -- Valor filtrado no tempo $t$
  \item $Y_t$ -- Leitura bruta instantânea do sensor
  \item $\alpha$ -- Fator de suavização ($\alpha = 0,1$ no sistema padrão)
\end{itemize}

O filtro é reinicializado automaticamente em cada novo ponto de ensaio para evitar o arraste de estados de pressão residuais de coletas anteriores.

\subsection*{1.2.2 Partição de Massa por Integral de Pressão Efetiva}
Em ensaios de pastas viscoplásticas, a massa total recolhida na balança ($M_{total}$) contém tanto o material extrudado durante o regime estacionário quanto o material expelido durante a fase transiente de pressurização. Para isolar a massa do regime estacionário ($m$), o software v4.1.1 utiliza o método da Integral de Pressão Efetiva ($I_{eff}$), baseando-se no fato de que a vazão é proporcional à pressão que excede a tensão de escoamento ($P_{yield}$):

\begin{equation}
    I_{eff} = \int_{0}^{t_{final}} \max(0, P(t) - P_{yield}) \, dt
\end{equation}

A fração da massa pertencente ao regime estacionário ($\text{Ratio}_{reg}$) é então calculada pela razão das integrais:

\begin{equation}
    \text{Ratio}_{reg} = \frac{\int_{t_{start}}^{t_{final}} \max(0, P(t) - P_{yield}) \, dt}{I_{eff}}
\end{equation}

A massa corrigida para o cálculo da taxa de cisalhamento é dada por $m = M_{total} \cdot \text{Ratio}_{reg}$. O sistema exibe um alerta visual caso $\text{Ratio}_{reg} < \SI{30}{\percent}$, indicando que o tempo em regime estacionário foi insuficiente frente à rampa de pressurização.

\subsection*{1.2.3 Cálculos Reológicos do Reômetro Capilar}
A análise reológica busca descrever a relação entre a tensão aplicada a um material e a taxa de deformação resultante. Os parâmetros fundamentais que descrevem essa relação são: a \textbf{tensão de cisalhamento} ($\tau$), que representa a força por unidade de área necessária para induzir o escoamento; a \textbf{taxa de cisalhamento} ($\dot{\gamma}$), que mede a taxa com que o material se deforma ou a velocidade relativa entre camadas adjacentes do fluido; e a \textbf{viscosidade} ($\eta$), definida como a razão entre a tensão e a taxa de cisalhamento ($\eta = \tau / \dot{\gamma}$), representando a resistência interna do material ao fluxo.

Para pastas e suspensões concentradas, a \textbf{tensão de escoamento} ($\tau_0$) é outro parâmetro importante, indicando a tensão mínima que deve ser superada para que o material comece a fluir.

No reômetro capilar, esses parâmetros são obtidos indiretamente a partir de gran\-de\-zas experimentais. A metodologia, aqui detalhada, baseia-se na medição da pressão de extrusão, da massa de material extrudado em um determinado tempo e nas dimensões conhecidas do capilar.

\textbf{Tensão de Cisalhamento na Parede:} O primeiro parâmetro calculado é a tensão de cisalhamento na parede ($\tau_w$). Sua equação é derivada de um balanço de forças em um elemento de volume cilíndrico de fluido em regime de escoamento permanente. A força motriz, gerada pela queda de pressão ao longo do capilar, é equilibrada pela força de atrito viscoso na parede. Isolando a tensão de cisalhamento, obtém-se a equação dada por \cite{barnes1989}:

\begin{equation}
    \tau_w = \frac{P \cdot R}{2 \cdot L}
\end{equation}

\textbf{Onde:}
\begin{itemize}[leftmargin=1.5em]
  \item $\tau_w$ -- Tensão de cisalhamento na parede do capilar (\si{Pa})
  \item $P$ -- Pressão de extrusão aplicada, medida pelo transdutor (\si{Pa})
  \item $R$ -- Raio interno do capilar (\si{m})
  \item $L$ -- Comprimento do capilar (\si{m})
\end{itemize}

Esta equação mostra que a tensão de cisalhamento é diretamente proporcional à pressão aplicada e ao raio do capilar, e inversamente proporcional ao seu comprimento.

\vspace{1em}
\textbf{Taxa de Cisalhamento Aparente:} O segundo parâmetro é a taxa de cisalhamento aparente na parede ($\dot{\gamma}_{aw}$). Ela é calculada a partir da vazão volumétrica ($Q$), que por sua vez é determinada pela massa extrudada, o tempo e a densidade da pasta. A equação para a taxa de cisalhamento aparente é dada por \cite{reed1995}:

\begin{equation}
    \dot{\gamma}_{aw} = \frac{4 \cdot Q}{\pi \cdot R^3}
\end{equation}

\textbf{Onde:}
\begin{itemize}[leftmargin=1.5em]
  \item $\dot{\gamma}_{aw}$ -- Taxa de cisalhamento aparente na parede (\si{s^{-1}})
  \item $Q$ -- Vazão volumétrica da pasta através do capilar (\si{m^3/s})
  \item $R$ -- Raio interno do capilar (\si{m})
  \item $\pi$ -- Constante matemática ($\approx$ 3,1416)
\end{itemize}

A vazão volumétrica $Q$ é calculada a partir das medidas experimentais diretas:
\begin{equation}
    Q = \frac{m}{\rho \cdot t}
\end{equation}

\textbf{Onde:}
\begin{itemize}[leftmargin=1.5em]
  \item $m$ -- Massa de pasta extrudada, medida pela balança (\si{kg})
  \item $\rho$ -- Densidade da pasta cerâmica (\si{kg/m^3})
  \item $t$ -- Tempo de extrusão (\si{s})
\end{itemize}

\textbf{Importante:} Ressalta-se o termo \textit{aparente}. Esta equação é derivada sob a premissa de um perfil de velocidade parabólico no interior do capilar, o que é estritamente válido apenas para fluidos Newtonianos. Como as suspensões cerâmicas são fluidos não-Newtonianos, seu perfil de velocidade não é parabólico, o que torna este valor uma aproximação inicial que será corrigida posteriormente.

\vspace{1em}
\textbf{Viscosidade Aparente:} A partir dos dois parâmetros anteriores, a viscosidade aparente ($\eta_a$) é calculada como a simples razão entre eles:

\begin{equation}
    \eta_a = \frac{\tau_w}{\dot{\gamma}_{aw}}
\end{equation}

\textbf{Onde:}
\begin{itemize}[leftmargin=1.5em]
  \item $\eta_a$ -- Viscosidade aparente do material (\si{Pa.s})
  \item $\tau_w$ -- Tensão de cisalhamento na parede (\si{Pa})
  \item $\dot{\gamma}_{aw}$ -- Taxa de cisalhamento aparente na parede (\si{s^{-1}})
\end{itemize}

Esta viscosidade também é chamada de \textit{aparente} porque é calculada usando valores não corrigidos da taxa de cisalhamento. Para fluidos não-Newtonianos como pastas cerâmicas, são necessárias correções adicionais para obter a viscosidade verdadeira.

\subsubsection*{Correções para Fluidos Não-Newtonianos}
Para uma análise rigorosa de fluidos não-Newtonianos, o uso de valores aparentes exige a aplicação de correções.

\textbf{Weissenberg-Rabinowitsch:} A mais importante delas para este estudo é a correção de Weissenberg-Rabinowitsch, que converte a taxa de cisalhamento aparente ($\dot{\gamma}_{aw}$) na taxa de cisalhamento verdadeira na parede ($\dot{\gamma}_w$). Esta correção é necessária porque, em fluidos pseudoplásticos, o perfil de velocidade do escoamento é mais achatado do que o perfil parabólico Newtoniano. A equação é dada por \cite{barnes1989}:

\begin{equation}
    \dot{\gamma}_w = \frac{\dot{\gamma}_{aw}}{4} \left( 3 + \frac{d(\ln \dot{\gamma}_{aw})}{d(\ln \tau_w)} \right) = \dot{\gamma}_{aw} \left( \frac{3n' + 1}{4n'} \right)
\end{equation}

Na versão 4.1.1, o software substituiu a aproximação de um índice pseudoplástico global ($n'$) pelo cálculo do perfil de derivadas locais ($\approx \Delta \ln \tau_w / \Delta \ln \dot{\gamma}_{aw}$) ponto a ponto, aplicando o fator adequado a cada faixa de cisalhamento para modelar com rigor fluidos que transitam entre comportamentos Newtonianos e viscoplásticos em diferentes gradientes.

\subsubsection*{Modelos de Ajuste Reológico}
O sistema avalia e ajusta automaticamente os principais modelos matemáticos utilizados na reologia de pastas cerâmicas:

\begin{enumerate}[label=(\alph*)]
  \item \textbf{Lei da Potência (Ostwald-de Waele):} Descreve fluidos pseudoplásticos sem tensão de escoamento.
  \begin{equation} \tau = K \cdot \dot{\gamma}^n \end{equation}
  \item \textbf{Bingham:} Para fluidos com comportamento linear após superar $\tau_0$.
  \begin{equation} \tau = \tau_0 + \eta_p \cdot \dot{\gamma} \end{equation}
  \item \textbf{Herschel-Bulkley:} Combina o modelo de Lei da Potência com tensão de escoamento.
  \begin{equation} \tau = \tau_0 + K \cdot \dot{\gamma}^n \end{equation}
  \item \textbf{Casson:} Comumente utilizado para pastas com escoamento plástico.
  \begin{equation} \sqrt{\tau} = \sqrt{\tau_0} + \sqrt{\eta_c \cdot \dot{\gamma}} \end{equation}
\end{enumerate}

\subsubsection*{Otimização e Criteração de Seleção (WLS e AIC)}
Historicamente utilizando a regressão de Mínimos Quadrados Ordinários (OLS) e escolhendo modelos com maior $R^2$, o sistema agora adota Mínimos Quadrados Ponderados (WLS, do inglês \textit{Weighted Least Squares}), minimizando a soma ponderada dos resíduos:

\begin{equation}
    S = \sum_{i=1}^n w_i (\tau_{i, med} - \tau_{i, pred})^2, \quad \text{com } w_i = \frac{1}{\sigma_i^2}
\end{equation}

Utilizando como base do peso ($w_i$) o desvio padrão instrumental propagado, confere-se às faixas com menor dispersão estatística maior influência na convergência final da curva de ajuste. A seleção do modelo ideal migrou da métrica de $R^2$ para o Critério de Informação de Akaike (AIC):

\begin{equation}
    AIC = 2k + n \ln\left(\frac{RSS}{n}\right)
\end{equation}

\textbf{Onde:}
\begin{itemize}[leftmargin=1.5em]
  \item $k$ -- Número de parâmetros do modelo (ex: 3 para Herschel-Bulkley)
  \item $n$ -- Número de pontos experimentais
  \item $RSS$ -- Soma dos quadrados dos resíduos (\textit{Residual Sum of Squares})
\end{itemize}

Este critério é fundamental para prevenir o \textit{overfitting} (superajuste), pois penaliza a adição de parâmetros extras que não reduzem significativamente o erro residual. O modelo vencedor é aquele que apresenta o menor valor de AIC.

\vspace{1em}
\textbf{Bagley:} Leva em conta as perdas de pressão na entrada e saída do capilar \cite{barnes1989}. A pressão total durante a extrusão capilar ($P$) inclui não apenas o esforço necessário para o escoamento dentro do capilar, mas também perdas associadas às extremidades ($P_e$):
\begin{equation}
P = 2 \tau_w \left( \frac{L}{R} \right) + P_e
\end{equation}

\textbf{Mooney:} Em suspensões concentradas, pode ocorrer deslizamento na parede (\textit{wall slip}) \cite{barnes1989}. A taxa de cisalhamento aparente é relacionada com a taxa de cisalhamento verdadeira do fluido ($\dot{\gamma}_{s,f}$) e a velocidade de deslizamento ($C$) pela equação:
\begin{equation}
\dot{\gamma}_{aw} = \dot{\gamma}_{s,f} + \frac{C}{R}
\end{equation}

\subsection*{1.3 Validação Experimental}
As versões anteriores do reômetro capilar foram validadas experimentalmente por meio de comparação direta com um reômetro comercial de referência (\textbf{Anton Paar MCR 102}), equipado com geometria de placas paralelas. Os ensaios comparativos, realizados com suspensões de caulim a \SI{60}{\percent} de sólidos, demonstraram boa concordância entre os equipamentos, especialmente para taxas de cisalhamento elevadas (acima de \SI{100}{s^{-1}}), confirmando a confiabilidade do dispositivo de baixo custo para análises reológicas.

\subsection*{1.4 Dicas Práticas de Preparação de Amostras}
Com base na experiência acumulada durante o desenvolvimento do equipamento, recomenda-se:

\begin{itemize}
  \item \textbf{Homogeneização:} Utilizar sacos plásticos resistentes para misturar caulim e água por agitação manual vigorosa. Este método facilita a homogeneização e permite transferência direta da pasta para o reômetro (cortando uma das extremidades do saco).
  
  \item \textbf{Limites de Concentração:} Para suspensões puras de caulim-água (sem aditivos), o limite prático de concentração é de aproximadamente \SI{65}{\percent} de sólidos. Acima deste valor, a pasta torna-se excessivamente rígida e quebradiça ("aspecto de farofa"), inviabilizando a extrusão sem o uso de modificadores reológicos (plastificantes ou defloculantes).
\end{itemize}

\section*{2. Variáveis e Nomenclatura Utilizada}
\begin{center}
\begin{tabular}{@{}l l@{}}
\toprule
\textbf{Símbolo} & \textbf{Descrição e Unidade} \\
\midrule
$R$ & Raio interno do capilar (\si{m}) \\
$D$ & Diâmetro interno do capilar (\si{mm}) \\
$L$ & Comprimento do capilar (\si{m}) \\
$Q$ & Vazão volumétrica (\si{m^{3}/s}) \\
$\rho$ & Densidade da pasta (\si{kg/m^{3}}) \\
$m$ & Massa extrudada (\si{kg}) \\
$t$ & Tempo de extrusão (\si{s}) \\
$P$ & Pressão aplicada (\si{Pa}) \\
$\tau_w$ & Tensão de cisalhamento na parede (\si{Pa}) \\
$\dot{\gamma}_{aw}$ & Taxa de cisalhamento aparente (\si{s^{-1}}) \\
$\dot{\gamma}_w$ & Taxa de cisalhamento real (corrigida) (\si{s^{-1}}) \\
$\eta_a$ & Viscosidade aparente (\si{Pa.s}) \\
$\eta$ & Viscosidade verdadeira (\si{Pa.s}) \\
$n'$ & Índice de pseudoplasticidade estimado (adimensional) \\
$K'$ & Coeficiente de consistência estimado (\si{Pa.s^n}) \\
$\tau_0$ & Tensão de escoamento (\si{Pa}) \\
$\eta_p$ & Viscosidade plástica (\si{Pa.s}) \\
$I_{eff}$ & Integral de pressão efetiva (\si{Pa.s}) \\
$\text{Ratio}_{reg}$ & Fração de massa em regime estacionário (adimensional) \\
$\alpha$ & Fator de suavização do filtro EMA (adimensional) \\
$AIC$ & Critério de Informação de Akaike (adimensional) \\
$RSS$ & Soma dos quadrados dos resíduos (\si{Pa^2}) \\
\bottomrule
\end{tabular}
\end{center}

\section*{3. Fluxo de Trabalho e Arquitetura do Sistema}

O sistema de análise reológica foi reestruturado na versão 4.1.1 para uma arquitetura modular baseada em pacotes, garantindo maior robustez e manutenibilidade. A operação principal agora é centrada na Interface Gráfica (GUI), enquanto os scripts de linha de comando originais foram preservados como opções de legado ou para uso "headless".

\noindent
\href{https://www.mermaidchart.com/d/6f9e496c-626c-429a-ba93-17c761bf78d0}{\textbf{Clique aqui para visualizar o Fluxograma do Processo}}

\subsection*{3.0 Interface Gráfica (gui\_main.py) -- V4.1.1}
A versão atual do sistema opera através de uma interface moderna desenvolvida com \textit{CustomTkinter}, organizada no pacote \texttt{gui/}. Este é o ponto de entrada padrão:

\begin{itemize}
  \item \textbf{Módulo Central:} O arquivo \texttt{gui\_main.py} atua como lançador, inicializando a aplicação definida em \texttt{gui/app.py}. A versão 4.1.1 inclui protocolos de segurança industrial (\texttt{WM\_DELETE\_WINDOW}) para garantir que a desconexão do hardware seja síncrona ao fechamento da janela.
  \item \textbf{Estrutura Modular:} As funcionalidades foram segregadas em módulos específicos em \texttt{gui/frames/} (Coleta, Histórico, Calibração, Análise, Correções). A versão atual inclui um gerenciador avançado de \textbf{Limpeza de Dados/Outliers} interativo e um sistema de alertas de qualidade que monitora a integridade da partição de massa via integral.
  \item \textbf{Abstração de Hardware:} A comunicação com o Arduino é gerenciada pela classe \texttt{ReometerController} (\texttt{reometer\_controller.py}). O sistema agora utiliza um ADC de 10 bits com \textit{oversampling} de 4 níveis para estabilidade de ruído branco.
  \item \textbf{Otimização de Banco de Dados:} A persistência utiliza SQLite (\texttt{reometria.db}). Na versão 4.1.1, foram implementadas otimizações de \textit{caching} de consultas (verificação de calibração a cada 60s) e resolução do problema de múltiplas requisições sequenciais (\textit{N+1 queries}) na tela de histórico, garantindo fluidez mesmo em bancos de dados com centenas de ensaios acumulados.
\end{itemize}

\textbf{Execução:} \texttt{python gui\_main.py}

\subsection*{3.1 Scripts Legados (CLI)}
Os scripts originais de linha de comando, numerados de 1 a 5, foram movidos para o diretório \texttt{legacy/}. Eles mantêm a lógica funcional original e servem como referência ou \textit{backend} para operações específicas.

\subsubsection*{3.1.1 Coleta e Controle (Script 1)}
Responsável pela interface direta com transdutores em modo texto. Gera arquivos JSON compatíveis com o novo sistema.

\subsubsection*{3.1.2 Análise Reológica (Script 2)}
Realiza o processamento matemático "bruto" dos arquivos JSON. As equações e modelos (Newton, Power Law, Bingham, Herschel-Bulkley, Casson) permanecem os mesmos utilizados pela GUI (agora encapsulados em \texttt{modelos\_reologicos.py} e \texttt{reologia\_fitting.py}).

\subsubsection*{3.1.3 Comparativo e Outros}
Os scripts de tratamento estatístico e comparativo analítico também foram preservados em \texttt{legacy/} para garantir reprodutibilidade de análises antigas.

\subsection*{3.2 Testes Automatizados}
A robustez matemática é garantida por uma suíte de testes unitários (\texttt{test\_calculos.py}) que valida as equações de correções e ajustes de modelos contra resultados conhecidos.

% --- Referências Bibliográficas (Manual) ---
\begin{thebibliography}{9}

\bibitem{barnes1989}
BARNES, H. A.; HUTTON, J. F.; WALTERS, K. \textbf{An Introduction to Rheology}. Amsterdam: Elsevier, 1989.

\bibitem{reed1995}
REED, J. S. \textbf{Principles of Ceramics Processing}. 2. ed. New York: John Wiley \& Sons, 1995.

\end{thebibliography}

\end{document}